\chapter{Introducción a NeoBilbao}

\section{Bienvenidos al mundo de Ur-Bare}

\neofigure{workers_marching_shift.jpg}{Los trabajadores marchando en silencio hacia sus turnos bajo el cielo de acero.}

\begin{quote}
«El aire en el exoplaneta Ur-Bare siempre sabe a metal, ozono y sangre. Fuera del perímetro de la colonia, las tormentas de polvo ferroso rojo devoran cualquier estructura que no esté forjada en titanio macizo. Dentro, el sirimiri ácido corroe las vigas y difumina las pantallas de neón sobre el asfalto. Bienvenidos a NeoBilbao».
\end{quote}

\textbf{NeoBilbao} es un juego de rol de mesa de ciencia ficción ciberpunk obrera y noir corporativo, ambientado en una colonia industrial remota fundada por la expedición vasca de Euskal-Astra a mediados del siglo XXI. Inspirado en la historia industrial de la margen izquierda del Nervión, la siderurgia pesada de Altos Hornos y la lucha sindical frente a megacorporaciones implacables, NeoBilbao sitúa a los jugadores no en el papel de mercenarios cromados o hackers millonarios, sino en la piel de la clase trabajadora: soldadores de plasma, médicos de guardia, celadores quemados, técnicos de mantenimiento y enlaces vecinales que intentan preservar su humanidad y cuidar de los suyos en una metrópolis que les exprime la vida.

\section{Las dos eras de juego}

Este manual unifica el trasfondo y el sistema de juego para permitir dirigir partidas en dos periodos históricos contrastados:

\begin{itemize}
  \item \lead{NeoBilbao 2084 -- La Margen (Noir Corporativo y Falsa Meritocracia):}  
  Dieciocho años antes del colapso total. La colonia aún conserva una clase media urbana hacinada en el Distrito 7 («La Margen»). La violencia de Orekak-Corp se ejerce de forma soterrada mediante policía política de paisano (Los Auditores), burocracia opresiva, deudas de oxígeno y el espejismo de la movilidad social a través del \textbf{Programa Azul de Excelencia Ciudadana}. La tecnología es táctil, analógica y ruidosa, con pantallas CRT de fósforo verde, disquetes magnéticos y cables coaxiales.
  
  \item \lead{NeoBilbao 2102 -- Los Zócalos (Ciberpunk Obrero y la Ría de Acero):}  
  Tras la devastadora Guerra del Plasma (2088--2094), el régimen corporativo ha sellado la colonia en dos mundos herméticos: la paradisíaca \textbf{Cúpula Dorada} para la oligarquía superior y el abismo brutalista de los \textbf{Zócalos Industriales} para los trabajadores (\emph{Langileak}). En este futuro, los mineros y fundidores soportan jornadas de dieciséis horas picando y refinando Burdin-Urdin, unidos a sus familias reubicadas en la zona alta mediante el implante neuronal del \textbf{Velo} y el solemne pacto del \textbf{Etxearen Zina} (el Juramento del Linaje).
\end{itemize}

\section{Aviso sin spoilers: Un manual para todos los jugadores}

\begin{CajaAmbientacion}{Dossier desclasificado}
Este volumen está expresamente redactado para ser leído tanto por \textbf{jugadores} como por \textbf{Directores de Juego}. Presenta la sociedad, la geografía, las facciones, la tecnología y el sistema de reglas sin desvelar las revelaciones confidenciales, giros de trama ni secretos de búnker de los dos módulos oficiales de aventura (*Etxearen Zina 2102* y *La Lotería Azul 2084*). 

Sirve como guía de consulta en mesa para crear cuadrillas de personajes, sumergirse en la atmósfera de Ur-Bare y diseñar campañas y misiones propias.
\end{CajaAmbientacion}

\section{Pilares temáticos de la ambientación}

\begin{enumerate}
  \item \lead{La solidaridad como trinchera (Auzolan):} En un entorno donde las corporaciones privatizan hasta el aire que respiras, el apoyo mutuo y la lealtad comunitaria son la única fuerza que impide la aniquilación moral de los personajes.
  \item \lead{Tecnología pesada y degradada:} Aquí no hay magia cibernética limpia. Los implantes duelen, supuran y requieren mantenimiento constante; los cables son gruesos; las herramientas pesan; y las averías son letales.
  \item \lead{El peso de los afectos:} Los personajes no luchan por créditos abstractos. Luchan por el vecino que les prestó un filtro de aire, por el compañero atrapado en la fundición o por la promesa hecha a sus hijos.
\end{enumerate}

\section{Herramientas de seguridad en la mesa}

NeoBilbao explora temas adultos e intensos: explotación laboral, angustia psicológica, hacinamiento, separación familiar y violencia corporativa. Para garantizar que la mesa disfrute de una experiencia emocionante y segura, recomendamos acordar previamente:

\begin{itemize}
  \item \lead{Líneas y Velos:} Definir colectivamente antes de jugar qué temas están totalmente fuera de la partida (líneas) y cuáles pueden ocurrir en la ficción pero sin describirse en detalle gráfico (velos).
  \item \lead{La Tarjeta X:} Cualquier participante en la mesa puede tocar la tarjeta o levantar la mano para pausar la narración de inmediato si una escena cruza un límite incómodo, sin necesidad de dar justificaciones.
  \item \lead{La Pausa de Auzolan:} Tomar breves pausas de descompresión tras escenas de alto estrés emocional para comprobar el bienestar de los jugadores.
\end{itemize}
